
\documentclass{article}
\usepackage{colortbl}
\usepackage{makecell}
\usepackage{multirow}
\usepackage{supertabular}

\begin{document}

\newcounter{utterance}

\centering \large Interaction Transcript for game `wordle\_withclue', experiment `medium\_frequency\_words\_clue\_no\_critic', episode 13 with refineonly{-}Qwen3.5{-}9B.
\vspace{24pt}

{ \footnotesize  \setcounter{utterance}{1}
\setlength{\tabcolsep}{0pt}
\begin{supertabular}{c@{$\;$}|p{.15\linewidth}@{}p{.15\linewidth}p{.15\linewidth}p{.15\linewidth}p{.15\linewidth}p{.15\linewidth}}
    \# & \multicolumn{2}{c}{Player} && \multicolumn{2}{c}{Game Master} \\
    \hline

    \theutterance \stepcounter{utterance}  
    & & & \multicolumn{4}{p{0.6\linewidth}}{
        \cellcolor[rgb]{0.9,0.9,0.9}{
            \makecell[{{p{\linewidth}}}]{
                \texttt{\tiny{[P1$\langle$GM]}}
                \texttt{You are a language wizard who likes to guess words by using the given rules.} \\
\\ 
\texttt{Welcome to Wordle! You have six attempts to guess the target word, a valid English word of five lowercase letters (a{-}z). Please use the tags "explanation:" and "guess:" to provide a concise explanation for each guess.} \\
\\ 
\texttt{To help you make an informed guess, you will receive a clue for the word, such as} \\
\texttt{clue: snowy white.} \\
\\ 
\texttt{Here is an example guess based on the clue:} \\
\texttt{explanation: In the fairy tail Snow White, the girl is killed because she eats a poisoned apple. And the word apple has 5 letters.} \\
\texttt{guess: apple} \\
\\ 
\texttt{After each guess, your answer will be validated, and you will receive feedback indicating which letters are correct (green), which letters are correct but in the wrong position (yellow), and which letters are incorrect (red). This feedback can be useful in determining which letters to include or exclude in your next guess.} \\
\\ 
\texttt{For example, the feedback for "apple" might be:} \\
\texttt{guess\_feedback: a<yellow> p<yellow> p<green> l<yellow> e<red>} \\
\\ 
\texttt{The explanation should contain details about how both the guess\_feedback and clue work together in tandem to formulate a new guess.} \\
\\ 
\texttt{Do not generate any other text.} \\
\\ 
\texttt{Let's begin with your first guess.} \\
\\ 
\texttt{clue: baby conveyance} \\
            }
        }
    }
     \\ \\

    \theutterance \stepcounter{utterance}  
    & & \multicolumn{4}{p{0.6\linewidth}}{
        \cellcolor[rgb]{0.95,0.95,0.95}{
            \makecell[{{p{\linewidth}}}]{
                \texttt{\tiny{[P1$\leftrightarrow$P]}}
                \texttt{{[}DRAFT{]}} \\
\texttt{explanation: The clue "baby conveyance" strongly suggests the word "stroller," which is a common term for a device used to transport a baby. The word has exactly five letters, fitting the Wordle constraints.} \\
\texttt{guess: stroller} \\
            }
        }
    }
    & \\ \\

    \theutterance \stepcounter{utterance}  
    & \multicolumn{4}{p{0.6\linewidth}}{
        \cellcolor[rgb]{0.561,0.737,0.561}{
            \makecell[{{p{\linewidth}}}]{
                \texttt{\tiny{[P1$\rangle$GM]}}
                \texttt{explanation: The clue "baby conveyance" suggests a stroller, but "stroller" has 8 letters, which exceeds the 5{-}letter limit required for Wordle. A valid 5{-}letter synonym for a baby conveyance is "pram" (4 letters) or "crib" (4 letters), but neither fits perfectly. However, "stroller" is often confused with "stroll" (6 letters). Let's reconsider: common 5{-}letter baby conveyances include "pram" (too short), "crib" (too short), "cot" (too short). Wait, "stroller" is definitely wrong length. What about "bassinet"? Too long. "Car seat"? Too long. Perhaps the word is "prams" (plural)? No, Wordle is singular. Let's think of other terms. "Stroller" is the most common, but it's 8 letters. Is there a 5{-}letter word? "Prig"? No. "Crib"? 4. "Cot"? 3. "Pram"? 4. "Push"? No. Maybe the clue implies "stroller" but the word is actually "stroll"? No. Let's try a different angle. What if the word is "prams"? No. How about "bass"? No. Wait, "stroller" is 8 letters. The only common 5{-}letter word for a baby conveyance might be "pram" if we consider British English, but it's 4 letters. Actually, "prams" is 5 letters but plural. Let's try "cribs"? 5 letters, plural. But Wordle is usually singular. What about "cot"? No. Maybe the word is "stroller" and I'm miscounting? S{-}T{-}R{-}O{-}L{-}L{-}E{-}R is 8. Okay, let's look for synonyms. "Buggy"? 5 letters. Yes, "buggy" is a common term for a baby conveyance, especially in British English. It fits the 5{-}letter constraint.} \\
\texttt{guess: buggy} \\
            }
        }
    }
    & & \\ \\

    \theutterance \stepcounter{utterance}  
    & & & \multicolumn{2}{p{0.3\linewidth}}{
        \cellcolor[rgb]{0.9,0.9,0.9}{
            \makecell[{{p{\linewidth}}}]{
                \texttt{\tiny{[GM$|$GM]}}
                \texttt{attempts = 1} \\
\texttt{target\_word = buggy} \\
\texttt{target\_word\_clue = baby conveyance} \\
\texttt{guess = buggy} \\
\texttt{guess\_feedback = b<green> u<green> g<green> g<green> y<green>} \\
            }
        }
    }
    & & \\ \\

    \theutterance \stepcounter{utterance}  
    & & & \multicolumn{2}{p{0.3\linewidth}}{
        \cellcolor[rgb]{0.9,0.9,0.9}{
            \makecell[{{p{\linewidth}}}]{
                \texttt{\tiny{[GM$|$GM]}}
                \texttt{game\_result = WIN} \\
            }
        }
    }
    & & \\ \\

\end{supertabular}
}

\end{document}
